
%(BEGIN_QUESTION)
% Copyright 2007, Tony R. Kuphaldt, released under the Creative Commons Attribution License (v 1.0)
% This means you may do almost anything with this work of mine, so long as you give me proper credit

The effects of {\it Lenz's Law} may be seen if we quickly pass a strong magnet near a thin strip of electrically conductive material, such as aluminum foil:

$$\includegraphics[width=15.5cm]{i02300x01.eps}$$

What will happen to the foil as the magnet is quickly swept past it?  Be as detailed as you can in your explanation.

\vskip 10pt

Now, imagine bending this strip of aluminum foil into a loop, and passing the same (powerful) magnet past the surface of the strip, except this time the magnet's motion must be circular because the strip's geometry is circular:

$$\includegraphics[width=15.5cm]{i02300x02.eps}$$

Explain how this ``thought experiment'' relates to the operation of an AC {\it induction} motor and the concept of a rotating magnetic field.  Also, explain why the rotor of an AC induction motor need merely be made of an electrically conductive material (such as copper or aluminum), not necessarily a magnetic material (such as iron).

\underbar{file i02300}
%(END_QUESTION)





%(BEGIN_ANSWER)

As a powerful magnet sweeps by a strip of aluminum foil, eddy currents are induced in the foil which produce their own reactionary magnetic fields.  In accordance with Lenz's Law, these magnetic fields produced by the eddy currents react against the powerful magnet's field in such a way that they oppose its relative change in position.  Thus, the aluminum foil strip will be ``dragged'' along by the field of the moving magnet.

\vskip 10pt

I'll let you determine what the aluminum foil loop will do!

%(END_ANSWER)





%(BEGIN_NOTES)

A very common misconception is that the rotor of an induction motor must be magnetic in nature in order to experience a torque from the rotating magnetic field.  Since Lenz's Law is a product of electromagnetic {\it induction}, the only criterion for the rotor material is that it be capable of supporting eddy currents (i.e. it must be a conductor of electricity).

That being said, most induction motor rotors contain a substantial amount of iron.  This is done to enhance the strength of the magnetic field from stator pole to stator pole (by reducing the magnetic reluctance between poles).

%INDEX% Electronics review: induction versus synchronous motor operation
%INDEX% Electronics review: Lenz's Law
%INDEX% Electronics review: synchronous versus induction motor operation

%(END_NOTES)


